\section{数据预处理}

\subsection{分析单位与重复测序合并}

男胎原始工作表包含1082条测序记录和267名孕妇。数据中既有同一孕妇的多次采血，也有同次采血的重复测序；部分重复结果虽然报告日期相差数日，但孕妇代码、检测抽血次数和记录孕周均相同。为避免将延后报告的技术复测误认为新的生物样本，本文统一以“孕妇代码、检测抽血次数”识别一次采血，对该次采血的Y染色体浓度、孕周、BMI和年龄取均值，并保留浓度最小值、最大值和测序次数用于误差分析。共识别40组同次采血重复测序，合并后得到1021次采血观测。

\begin{table}[H]
  \centering
  \caption{问题一、问题二统一的数据清洗统计}
  \label{tab:male-clean-summary}
  \begin{tabular}{cc}
    \toprule
    项目 & 数量 \\
    \midrule
    原始检测记录 & 1082条 \\
    男胎孕妇 & 267人 \\
    聚合后采血观测 & 1021次 \\
    同次采血重复测序组 & 40组 \\
    \bottomrule
  \end{tabular}
\end{table}

\subsection{变量标准化与达标状态构造}

将“$a\mathrm{w}+b$”形式的孕周转换为连续周数：

\begin{equation}
  W=a+\frac{b}{7}.
  \label{eq:week-convert}
\end{equation}

利用体重除以身高平方复核BMI，差异均小于0.02。问题一使用各次采血记录对应的BMI解释连续Y染色体浓度；为避免把达标之后发生的体重变化反向用于预测，问题二采用每名孕妇首次检测时的BMI作为基线BMI。定义第$i$名孕妇第$j$次采血的达标状态为

\begin{equation}
  D_{ij}=\mathbb{I}\!\left(Y_{ij}\geq 0.04\right),
  \label{eq:q2-status}
\end{equation}

其中$Y_{ij}$表示Y染色体浓度。GC含量略低于$40\%$的记录仅设置质量标记而不机械删除，避免一次性剔除大量临界样本。

\subsection{删失结构说明}

267名孕妇中，217人在第一次检测时已经达标，只能确定真实达标时间早于首次检测，属于左删失；43人的真实达标时刻位于相邻两次检测之间，属于区间删失；7人在最后一次检测时仍未达标，属于右删失。因此，首次观测到达标的孕周不能直接等同于精确的生理达标时间，后续模型应尽量使用全部纵向检测记录。
